% Verbatim recovery from nucleo.tex REV08, lines 10--32.
% Extension specific to II, placed after the unchanged common kernel.
\subsection{Correspondence between positional support and the oriented calendar}
\label{ii:ampliacion-posicional-calendario}

The calendar representation distinguishes the positions of APP from the
oriented transitions of its evolution. We revisit this correspondence to
make explicit the reading of the segment containing the junction between turns.

Consider the alphabet of nine indices
\[
 D_9=\{1,2,3,4,5,6,7,8,9\}.
\]
Its horizontal and vertical arrangement produces the eighty-one positions
of \(D_9^2\). The oriented cycle \(C_9\) is also used to describe
motion. Its nine edges are assigned, in order, the indices of
\(D_9\). The index of an edge specifies a transition between states,
whereas a vertex specifies its boundary. The map assigning
each edge its index allows the same alphabet to be used in the positional
support and in the calendar; it preserves the distinction between position,
oriented interval, and transported state.

In the integer lift of the calendar, the first turn comprises
the nine consecutive intervals \([0,1],\ldots,[8,9]\). The continuation
\([9,10]\) starts the next turn with index one and updated
memory. The segment drawn up to ten shows both turns at their
junction. This representation specifies the nine transitions of
one turn and the following transition; the holonomy of the state is constructed
subsequently through the effective composition of the TPK.

\begin{figure}[htbp]
\centering
\begin{tikzpicture}[x=1.05cm,y=1cm,>=Latex,font=\small]
 \draw[very thick,ink] (0,0)--(9,0);
 \draw[very thick,dashed,accent] (9,0)--(10,0);
 \foreach \x in {0,...,10}{
   \draw (\x,-.08)--(\x,.08);
   \node[below=5pt,font=\scriptsize] at (\x,0) {$\x$};
 }
 \foreach \x in {1,...,9}{
  \draw[->,ink] (\x-.85,.35)--(\x-.15,.35);
  \node[above=3pt] at (\x-.5,.35) {$\x$};
 }
 \draw[->,accent] (9.15,.35)--(9.85,.35);
 \node[above=3pt,text=accent] at (9.5,.35) {$1$};
 \draw[decorate,decoration={brace,amplitude=5pt},ink] (0,1)--(9,1)
  node[midway,above=7pt] {Nine oriented intervals: one turn};
 \node[align=center] at (5,-1.05)
  {Transition index: $g(t)=1+(t\bmod9)$\\
   Memory of the lifted calendar: $m(t)=\lfloor t/9\rfloor$};
 \node[align=center] at (5,-1.85)
  {$g(t+9)=g(t),\qquad m(t+9)=m(t)+1$};
\end{tikzpicture}
\caption{Lift of the cycle of nine transitions. The upper numbers
label oriented intervals; the lower numbers locate their boundaries.
The dashed interval shows the first transition of the next turn.
The phase index returns and the calendar quotient increases. This diagram
represents the calendar; the action on residues, quotients, orientation,
boundary, and memory is defined through the operations of the TPK.}
\label{fig:ii-segmento-generador}
\end{figure}

